\documentclass[10pt,a4paper]{amsart}
\usepackage[margin=19mm]{geometry}
\usepackage{amsmath,amssymb,mathtools}
\usepackage[scr=rsfs,cal=euler]{mathalfa}
\usepackage{xfrac}
\usepackage[unicode,hidelinks,bookmarks=false]{hyperref}
\newcommand{\ie}{i.e.\ }
\newcommand{\cf}{cf.\ }
\newcommand{\resp}{respectively\ }
\newcommand{\Subsection}{\subsection}
\providecommand{\nobreakdash}{\nobreak\hbox{-}}
\setlength{\emergencystretch}{2em}
\multlinegap0pt
\usepackage{amssymb}
\usepackage{float}
\usepackage{tikz}
\newtheorem{lemma}{Lemma}
\newcommand\ale[1]{\marginpar{Alejandro: #1}}
\newcommand\cat[1]{{}^\curvearrowright #1}
\DeclareMathOperator{\dom}{dom}
\def\s{\subseteq}
\newcommand\tom[1]{\marginpar{Tom: #1}}
\title{Compactness phenomena in HOD and the Optimality of Magidor's Covering theorem}
\author{Tom Benhamou, James Cummings, Yair Hayut, Gabriel Goldberg, Alejandro Poveda}
\date{Source: arXiv:2608.24190v1 (CC BY 4.0)}
\begin{document}
\maketitle

\noindent\emph{Source and licence.} Compactness phenomena in HOD and the Optimality of Magidor's Covering theorem. Version arXiv:2608.24190v1 (CC BY 4.0), distributed by arXiv under the Creative Commons Attribution 4.0 International licence, http://creativecommons.org/licenses/by/4.0/, as recorded in the arXiv licence metadata for this exact version and shown on the abstract page https://arxiv.org/abs/2608.24190v1. Full licence notice: \texttt{licenses/arxiv-2608.24190-CC-BY-4.0.txt}.\par\smallskip

\noindent\textit{Editorial scope.} Counted unit: the article's lemma residue (source0.tex lines 1045--1049). The article's complete original proof of it is delivered with it (source0.tex lines 1050--1079, one \texttt{proof} environment of 30 lines). No mathematics is written, repaired or re-derived: the statement and the proof are the article's own lines, in the article's own notation, and no retained line is substituted or re-flowed.  No further source-local context is needed: every cross-reference of every retained line resolves inside the two delivered spans.  Nothing was omitted from any retained span, so the package carries no omission record.  No retained line contains a citation, so the wrapper carries no bibliography and no bibliography entry.  No retained line contains a figure environment, an image-inclusion command or drawing code, so no image is delivered and the package declares no asset dependency.  Editorial formatting only, in addition: an AMS \texttt{amsart} preamble that restates the article's own declarations and macros used by the retained lines, namely the environments \texttt{lemma} on separate counters, together with the standard packages those lines need.  The retained environments print in this document as Lemma 1 in order of appearance, whereas the article prints the same items with the numbering of its own full document.  The complete native source of the article as distributed in the arXiv e-print for this exact version is bundled unchanged as \texttt{sources/208455/source0.tex}.
\begin{lemma}\label{lemma: residue}
    There is a  residue map $\restriction_M\colon \mathbb{P}\to \mathbb{P}\cap M$. %\tom{what is the difference between residue map and a projection? was this defined somewhere?}
   % \ale{The notion is due to Mitchell –– it appears on work of his on the $I[\omega_2]$. Strong residue maps need not be order-preserving.}\tom{I think we need to define it properly in the preliminaries}$\restriction_M\mathbb{P}\to \mathbb{P}_M$.
% In particular, the latter is a complete subposet of the former.
\end{lemma}

\begin{proof}
Let us define $ \restriction_M$ as follows: Given a condition $p=p_{\leftarrow}{}^\smallfrown \langle f, T\rangle$ in $\mathbb{P}$, $$p\restriction_M:=p_{\leftarrow}{}^\smallfrown \langle f\restriction M, T\restriction M\rangle,$$
where $T\restriction M$ is defined as  $\{\langle \nu_0\restriction M,\dots, \nu_{n-1}\restriction M\rangle\mid \langle \nu_0,\dots, \nu_{n-1}\rangle\in T\}.$

The first observation is that $p\restriction_M$ is a condition in $\mathbb{P}_M$. This is a routine checking using the fact that $V_\kappa\s M$ and that ${}^\kappa M\s M$. 

\smallskip

Let us show that $\restriction_M$ is a residue map. Let $r\leq_{\mathbb{P}_M} p\restriction_M$, $r\in M$, and let us show that $r$ is $\leq_{\mathbb{P}}$-compatible with $p$. First, by elementarity, $r\leq_{\mathbb{P}_M} p\restriction_M$ is the same as $r\leq_{\mathbb{P}} p\restriction_M$. By definition, there is $\vec\nu\in T\restriction M$ such that $$\text{$r_\leftarrow \leq (p\restriction_M\cat \vec\nu)_{\leftarrow}$ and $r_\rightarrow\leq^*_{\mathbb{P}} (p\restriction_M\cat \vec\nu)_\rightarrow$}.\footnote{Here $\leq$ denotes the product ordering associated to the `lower' extender based forcing.}$$
To simplify notations let us assume that $\vec\nu =\langle \nu\rangle$.

Then, $r_\leftarrow$ decomposes as $r_{\leftarrow\leftarrow}{}^\smallfrown r_{\leftarrow\rightarrow}$ and similarly $(p\restriction_M\cat \langle \nu\rangle )_{\leftarrow}$ decomposes as $(p\restriction_M\cat \langle \nu\rangle)_{\leftarrow\leftarrow}{}^\smallfrown (p\restriction_M\cat \langle \nu\rangle)_{\leftarrow\rightarrow}$. Notice that $(p\restriction_M\cat \langle \nu\rangle)_{\leftarrow\leftarrow}$ is just $p_{\leftarrow}$ due to the definition of the map $\restriction_M$. Set $p'_0$ to be $r_{\leftarrow\leftarrow}$.

Now, let us analyze $(p\restriction_M\cat \langle \nu\rangle)_{\leftarrow\rightarrow}$. By definition of $\cat \langle \nu\rangle$, this is  $$\langle (f\restriction M)\downarrow \langle \nu\rangle, (T\restriction M)\downarrow \langle \nu\rangle\rangle;$$
more explicitly:
\begin{itemize}
    \item $(f\restriction M)\downarrow \langle \nu\rangle$ is the function with domain $\{\nu(\bar \alpha)\mid o(\bar\alpha)\neq 0\}$ and
    $$\left((f\restriction M)\downarrow \langle \nu\rangle\right)(\nu(\bar\alpha)):=(f\restriction M)(\bar \alpha)\restriction k_\nu=f(\bar \alpha)\restriction k_\nu,$$
    where $k_\nu:=\min\{l<|f(\bar\alpha)|\mid\forall l\leq i<|f(\bar\alpha)|\, (o(f_i(\bar\alpha))<o(\nu(\bar\alpha))) \}$.
    \item $(T\restriction M)\downarrow \langle \nu\rangle$ is the projected tree $$\{\langle \mu_0\downarrow \nu, \dots, \mu_{n-1}\downarrow\nu\rangle\mid \vec\mu\in (T\restriction M)\},$$
    which is the same as $\{\langle \mu_0\downarrow \nu, \dots, \mu_{n-1}\downarrow\nu\rangle\mid \vec\mu\in T\}.$
\end{itemize}
Let $\eta\in T$ be a $\dom(f)$-object such that $\eta\restriction M=\nu$. 

We define $p'_1$ to be $\langle f^{r_{\leftarrow\rightarrow}}\cup (f\downarrow \eta), S_1\rangle$ where $S_1$ is the intersection of the pullbacks to  $\dom(f^{r_{\leftarrow\rightarrow}})\cup \dom(f\downarrow\eta)$ of the trees $T^{r_{\leftarrow\rightarrow}}$ and $(T\downarrow\eta)$. 

Finally, let $p'_2$ to be $\langle f^{r_{\rightarrow}}\cup (f\cat \eta), S_2\rangle$ where $S_2$ is defined analogously.

Finally, we set $p'=p'_0{}^\smallfrown p'_1{}^\smallfrown p'_2$. It follows that this is a condition in $\mathbb{P}$ and that $p'$ witnesses compatibility of $r$ and $p$.
\end{proof}

\end{document}
